% ============================================================================
\chapter*{Эпилог. Дорожная карта замкнута: от уравнения самозамкнутости к единому Естествознанию}
\addcontentsline{toc}{chapter}{Эпилог. Дорожная карта замкнута}

Оглянемся на весь пройденный путь с первой до седьмой ступени:
\begin{enumerate}[label=\textbf{\arabic*.},leftmargin=*,noitemsep]
\item Из одного условия самозамкнутости конечного автомата без внешней памяти $\mathbf{1 + 2^d = d^2}$ мы вывели трёхмерность пространства $\mathbf{d = 3}$, арифметическую единицу $\mathbf{1 = 3^2 - 2^3}$, 9-битное фазовое кольцо $\mathcal{M} = \mathbf{\mathbb{Z}_{512}}$ и 12-лучевой кубооктаэдр $\mathbf{O_2(3)}$.
\item Мы увидели, что двусторонний гомеоморфизм носителей $\langle\mathcal{M},\mathcal{T};\,\pi_{M\to T},\,\iota_{T\to M}\rangle$ (в смысле теории моделей А.\,И.~Мальцева) сводит все естественные науки и все шесть отраслей инженерии к \textbf{трём законам классической механики Ньютона в точных целых числах на решётке $O_2(3)\times\mathbb{Z}_{512}$} и пяти механическим следствиям дискретности решётки ($c_0=\sqrt{2}c$, $D(\mathbf{x})\in[\varepsilon,833\varepsilon]$ с гало $16/3$, люфт зацепления $\Delta\varphi_{\mathrm{disc}}=41$, кольцо $\mathbb{Z}_{512}$ с упором Паули $256$ и 12-вальная оболочка $O_2(3)$ с седловиной $Q_\square=128$).
\item Каждая теорема этого курса от начала до конца формально проверена компилятором \textbf{Lean~4} (\texttt{0 sorryAx}, \texttt{0 Classical.choice}) и воплощена в открытом кремниевом RTL-ядре (\texttt{hardware/cuboctahedron\_carrier\_asic.v}), аналоговой 3D LC-решётке и лабораторных протоколах, доступных каждому исследователю и инженеру.
\end{enumerate}

\bigskip
\begin{center}
  \rule{0.65\textwidth}{0.4pt}\\[6pt]
  {\large\bfseries Дорожная карта замкнута: $1 + 2^d = d^2 \;\Longrightarrow\; \mathbb{Z}_{512} \;\Longrightarrow\; O_2(3) \;\Longrightarrow\; \text{Ньютон в }\mathbb{Z},\ \text{Кремний, Жизнь и Инженерия}$.}
\end{center}
